\section{Preparing an end with zero momentum}\label{sec:end}

Put $k=n-1$ and $p=n-2=k-1$.  A radius function $\varrho$ below
is smooth and positive on the closed exterior and equals the Euclidean
radius on a sufficiently distant part of its coordinate end.  Constants
in the end-preparation argument may depend on the original data and on its fixed
future-timelike ADM vector.  In particular, no uniform estimate as
$E-|P|\downarrow0$ is asserted. The final subsection supplies the positive-shell
construction used in
Section~\ref{sec:comparison} to prove automatic timelikeness; that
construction does not assume timelikeness.

The target class is Definition~\ref{def:prepared}. The present
task is to reach it while comparing every full cut with an original cut.

\begin{proposition}[Reduction to a prepared exterior]
\label{prop:end:reduction}
Let $(\Omega,g,K)$ be a smooth one-ended exterior as in
Definition~\ref{def:data}, with compact complement of its coordinate
end and with
\[
 g-\delta=O_6(r^{-q}),\qquad K=O_5(r^{-1-q}),\qquad
 \frac{n-2}{2}<q<n-2.
\]
Assume integrable constraint densities satisfying $\mu\ge|J|_g$,
finite ADM limits satisfying $E>|P|$, and $\theta_+(S)\le0$ on every
boundary component.  Suppose that the full-cut infimum of
Definition~\ref{def:fullcuts} is $a=A_{\min}^{g}(S)>0$, and put
$m=(E^2-|P|^2)^{1/2}$.
There are prepared data $(g_j,K_j)$ on the same exterior, with one
fixed $\beta\in(0,1)$, numbers $\epsilon_j\downarrow0$, and proper
diffeomorphisms $\Phi_j:\Omega\longrightarrow\Omega$ equal to the
identity near $S$, such that
\begin{equation}\label{eq:end:reduction-limits}
 P_j=0,\qquad E_j\longrightarrow m,\qquad
 \Phi_j^*g\le(1+\epsilon_j)g_j,\qquad
 A_{\min}^{g_j}(S)\ge(1+\epsilon_j)^{-k/2}a.
\end{equation}
The positive constant $c_j$ in Equation~\eqref{eq:end:prepared} and
the end coordinates may depend on $j$.
\end{proposition}

We first match the original energy and momentum with a
boosted vacuum end.  A compact annular correction joins the two ends
with a constraint error whose total scale tends to zero.  Within the
vacuum region we then bend the slice to a static one, giving zero
momentum.  A positive conformal factor repairs the small annular
error without decreasing the metric.  A proper radial comparison
map protects every full enclosing cut during the change of slice.
For each already fixed replacement, a second scalar perturbation makes
the energy condition and future trapping strict while changing the energy
by a quantity tending to zero. The argument fixes the original boost
throughout; it needs
no estimates uniform as the ADM vector approaches the null cone.

An asymptotic change to a rest frame is also used by Allen, Bryden,
Kazaras and Khuri in the invariant-mass step of
\cite[Section~5]{AllenBrydenKazarasKhuri2025}. Here the end is first
replaced by an explicit vacuum model, and the proper comparison map in
Proposition~\ref{prop:end:reduction} controls every full cut during
the replacement and bending.

\subsection{Conformal changes and a vacuum model with matching charges}

The conformal identities below explain how a scalar function can
improve the constraint inequality and the boundary expansion at the
same time.  We then compute the vacuum model's charges directly in
the ADM normalization used here.

For a smooth function $s$ set $g_s=e^{2s}g$ and $K_s=e^sK$.
If $|v|_g\le1$, the corresponding vector in the new unit ball is
$v_s=e^{-s}v$.  The conformal connection formula and scalar curvature
formula give
\begin{equation}\label{eq:end:conformal}
\begin{aligned}
 e^{2s}\bigl(\mu_s+J_s(v_s)\bigr)
 &=\mu+J(v)-k\Delta_gs-\frac{kp}{2}|ds|_g^2
                  +kK(\nabla s,v),\\
 e^s\theta_{+,s}&=\theta_++k\partial_\nu s.
\end{aligned}
\end{equation}
Here the normal $\nu$ points into the exterior.  For completeness,
$\tau_s=e^{-s}\tau$, $|K_s|_{g_s}^2=e^{-2s}|K|_g^2$, and
\[
 R_{g_s}=e^{-2s}\bigl(R_g-2k\Delta_gs-kp|ds|_g^2\bigr).
\]
Writing $\pi=K-\tau g$, contraction of
$\Gamma(g_s)^l_{ij}-\Gamma(g)^l_{ij}
=s_i\delta_j^l+s_j\delta_i^l-g_{ij}s^l$ yields
$J_{s,i}=e^{-s}(J_i+kK_{ij}\nabla^js)$.
The tangential trace of $K$ scales by $e^{-s}$ and
$H_s=e^{-s}(H+k\partial_\nu s)$, proving both identities.
In particular, if $s=2p^{-1}\log U$, $U>0$, the entire added term
in the first identity is
\begin{equation}\label{eq:end:conformal-U}
 \frac{2k}{pU}\bigl(-\Delta_gU+K(\nabla U,v)\bigr).
\end{equation}

\begin{lemma}[Charges of a boosted Schwarzschild--Tangherlini end]
\label{lem:end:boost}
In the Schwarzschild--Tangherlini spacetime of geometric mass $M>0$,
let $b,y$ be static time and isotropic spatial coordinates.  Make the
Lorentz change
\begin{equation}\label{eq:end:boost-coordinates}
 b=\gamma(T+vx^1),\qquad y^1=\gamma(x^1+vT),\qquad
 y^a=x^a\ (2\le a\le n),\qquad \gamma=(1-v^2)^{-1/2},
\end{equation}
where $0\le v<1$.  The induced data on $T=0$, oriented by the
future normal, have
\begin{equation}\label{eq:end:boost-charges}
 E=\gamma M,\qquad P_1=\gamma Mv,\qquad P_a=0\ (a>1).
\end{equation}
The metric and second fundamental form have the homogeneous leading
terms computed below, with remainders $O_d(r^{-2p})$ and
$O_d(r^{-2p-1})$, respectively, for every fixed derivative order $d$.
\end{lemma}

\begin{proof}
The exact isotropic metric is
\[
 -\left(\frac{1-M/(2|y|^p)}{1+M/(2|y|^p)}\right)^2db^2
 +\left(1+\frac{M}{2|y|^p}\right)^{4/p}|dy|^2.
\]
Its leading perturbation of Minkowski space is
$2M|y|^{-p}(db^2+p^{-1}|dy|^2)$.  On $T=0$ put
\[
 \rho^2=\gamma^2(x^1)^2+\sum_{a=2}^n(x^a)^2,\qquad f=\rho^{-p},
 \qquad A=\gamma^2(v^2+p^{-1}),\quad B=p^{-1},\quad
 C=kp^{-1}\gamma^2v.
\]
The nonzero relevant leading components of the spacetime perturbation
$\mathsf H$ are
$\mathsf H_{11}=2MAf$, $\mathsf H_{aa}=2MBf$, and
$\mathsf H_{01}=2MCf$.  Moreover
\[
 \partial_1f=-p\gamma^2x^1\rho^{-n},\qquad
 \partial_af=-px^a\rho^{-n},\qquad
 \partial_Tf=-p\gamma^2vx^1\rho^{-n}.
\]
Thus the energy flux vector $Q_i=\partial_j\mathsf H_{ij}
-\partial_i\mathsf H_{jj}$ satisfies
\[
 Q_1=-2MkB\partial_1f=2Mk\gamma^2x^1\rho^{-n},\qquad
 Q_a=-2M(A+pB)\partial_af=2Mp(A+1)x^a\rho^{-n}.
\]
Since $p(A+1)=k\gamma^2$, these expressions give the pointwise identity
\begin{equation}\label{eq:end:boost-energy-vector}
 Q_i=2Mk\gamma^2x_i\rho^{-n}.
\end{equation}

The chosen sign of the second fundamental form gives, to first order,
\[
 K_{ij}=\tfrac12\bigl(\partial_T\mathsf H_{ij}
                 -\partial_i\mathsf H_{0j}
                 -\partial_j\mathsf H_{0i}\bigr).
\]
Direct substitution gives
\begin{align*}
 K_{11}&=M\gamma^4v(2p+1-pv^2)x^1\rho^{-n},&
 K_{aa}&=-M\gamma^2vx^1\rho^{-n},\\
 K_{1a}&=Mk\gamma^2vx^a\rho^{-n},&
 K_{ab}&=0\quad(a\ne b).
\end{align*}
Their Euclidean trace is
$M\gamma^4v(p+v^2)x^1\rho^{-n}$.  Consequently, for the leading
Euclidean trace reversal $\pi_{ij}=K_{ij}-(\tr_\delta K)\delta_{ij}$,
\begin{equation}\label{eq:end:boost-momentum-tensor}
 \pi_{1j}=Mk\gamma^2v x_j\rho^{-n},\qquad
 \pi_{a1}=Mk\gamma^2vx^a\rho^{-n},\qquad
 \pi_{ab}=-Mk\gamma^4vx^1\rho^{-n}\delta_{ab}.
\end{equation}
Using the exact normal and the exact metric in the trace reversal
changes these formulas by $O_d(r^{-2p-1})$.  The energy remainder
has first derivatives of this same order.  Their sphere integrals
are $O(r^{-p})$, and therefore vanish for every $n\ge3$.

The ellipsoid $\{|\operatorname{diag}(\gamma,1,\ldots,1)x|<1\}$
has volume $\omega/(n\gamma)$.  Polar integration therefore gives
\begin{equation}\label{eq:end:ellipsoid-integral}
 \int_{S^{n-1}}|\operatorname{diag}(\gamma,1,\ldots,1)\theta|^{-n}
 \,d\theta=\frac{\omega}{\gamma}.
\end{equation}
Equations~\eqref{eq:end:boost-energy-vector} and
\eqref{eq:end:boost-momentum-tensor}, divided by $2k\omega$ and
$k\omega$, give $E=M\gamma$ and $P_1=M\gamma v$.
The integrand of a transverse momentum row is a constant times
$x^1x^a\rho^{-n}/r$, whose integral vanishes by reflection.
\end{proof}

For the original charges we now fix $M=m$, $v=|P|/E$, and rotate
the first axis to the direction of $P$ when $P\ne0$.
Lemma~\ref{lem:end:boost} gives a model end $(g^{\mathrm b},K^{\mathrm b})$
with exactly the same charges.  The parameter $v$ is fixed in all
limits that follow.

\subsection{Joining the original and model ends}

A direct cutoff need not preserve the dominant energy condition.
Our immediate task is to remove its leading constraint error by
compactly supported corrections.  The obstructions to a compact
correction are finitely many moments.  Charge matching removes their
constant limits, and integrability makes the remaining first moments
small enough after rescaling.

Localized constraint deformation and the finite-dimensional obstruction
from the adjoint kernel are central to the gluing methods of
Corvino--Schoen \cite{CorvinoSchoen2006} and Chru\'sciel--Delay
\cite{ChruscielDelay2003}. We use this methodological starting point,
not an imported rest-end replacement theorem: the present data need
not be vacuum, no parity is imposed, and both the energy inequality
and every full enclosing cut must survive the replacement.

We give the compact correction explicitly because neither parity
conditions nor convergence of center of mass or angular momentum
is assumed.  On Euclidean tensors define
\[
 \mathsf S(h)=\partial_i\partial_j(h_{ij}-(\tr_\delta h)\delta_{ij}),
 \qquad
 \mathsf D(b)_i=\partial_j(b_{ij}-(\tr_\delta b)\delta_{ij}).
\]
Their formal adjoint kernels relevant to compact support are,
respectively, affine scalar functions and Euclidean Killing fields.

\begin{lemma}[Compact right inverses]\label{lem:end:compact-inverses}
Let $\mathcal A\subset\R^n$ be a bounded connected Euclidean annulus
and fix a compact subset of its interior.
Smooth scalar sources supported in that subset with zero constant
and first moments have smooth symmetric solutions of
$\mathsf S(h)=f$ supported in a larger fixed compact subset of
$\mathcal A$.  Smooth vector sources with zero integrals against
every Euclidean Killing field have similarly supported symmetric
solutions of $\mathsf D(b)=F$.  For $1<a<\infty$ and integers $d\ge0$,
these solutions may be chosen linearly with estimates
\[
 \|h\|_{W^{d+2,a}}\le C_{d,a}\|f\|_{W^{d,a}},\qquad
 \|b\|_{W^{d+1,a}}\le C_{d,a}\|F\|_{W^{d,a}}.
\]
For arbitrary sources the same conclusions hold after subtracting
fixed smooth moment bumps, whose coefficients are bounded by the
absolute values of the respective moments.
\end{lemma}

\begin{proof}
We use the compact scalar divergence inverse on a connected open
set: a mean-zero smooth function supported in a fixed compact set
has a compactly supported vector primitive, with one Sobolev
derivative gained; the support and Sobolev mapping statements are
given, for example, by the regularized Bogovskii operator in
\cite[Section~3.1, Equation~(3.13), Theorem~3.2,
Corollaries~3.3--3.4, Remark~3.5,
and Proposition~4.1(ii)]{CostabelMcIntosh2010}.
Here is its local construction and the support
issue.  On a ball, choose a smooth unit-integral function supported
in a smaller concentric ball and use the Bogovskii integral
\[
 (\mathcal Bf)(x)=\int f(y)(x-y)
          \int_1^\infty\eta(y+t(x-y))t^{n-1}\,dt\,dy.
\]
Differentiation in distributions gives
$\divg\mathcal Bf=f-\eta\int f=f$; the segments occurring in the
integral remain in the ball.  For compact input support and compact
$\supp\eta$, their union is compact in that ball.  The kernel is of
order $1-n$, and its differentiated kernel has the cancellation
of a Calderon--Zygmund kernel.  The resulting $W^{d,a}$ estimate
is the usual boundedness of that singular integral and its
commutators with derivatives.  Cover the prescribed support by
finitely many balls compact in $\mathcal A$, connect them by a
finite chain of overlapping such balls, and use a partition of
unity.  Transfer the means along the chain using unit-integral
bumps in the overlaps.  Each resulting ball source has mean zero;
the sum of its ball primitives proves the stated inverse on
$\mathcal A$, with fixed support and constants.

Apply it first to a scalar $f$ with its affine moments zero,
obtaining $\partial_iv_i=f$.  Compact support gives
$\int v_i=-\int z_i f=0$.  Apply the same inverse to each $v_i$
to obtain $\partial_jT_{ij}=v_i$.  Then
$Q=(T+T^{\mathsf t})/2$ has $\partial_i\partial_jQ_{ij}=f$.
The tensor $h=Q-k^{-1}(\tr_\delta Q)\delta$ has trace reversal
$Q$, so $\mathsf S(h)=f$ with the asserted two-derivative gain.

For a vector source $F$, first solve $\partial_jB_{ij}=F_i$
row by row.  Translation moments permit this solve.  Rotation
moments give
\[
 0=\int(z_iF_j-z_jF_i)=\int(B_{ij}-B_{ji}).
\]
Put $A_{ij}=(B_{ij}-B_{ji})/2$ and solve
$\partial_lD_{ijl}=-A_{ij}$, choosing $D_{ijl}=-D_{jil}$.
Set
\[
 E_{ij}=\partial_l(D_{ijl}-D_{ilj}-D_{jli}),\qquad C_{ij}=B_{ij}+E_{ij}.
\]
Index interchange gives $E_{ij}-E_{ji}=-2A_{ij}$ and
$\partial_jE_{ij}=0$: the first two double derivatives cancel
after interchanging $j,l$, and the last vanishes by skew symmetry
in those indices.  Thus $C$ is symmetric and $\partial_jC_{ij}=F_i$.
Taking $b=C-k^{-1}(\tr_\delta C)\delta$ proves the assertion,
including the derivative count.

Finally take a nonnegative smooth bump supported in an open ball
of $\mathcal A$.  The Gram matrix of the affine functions, or
of the vector Killing fields, with this weight is positive
definite: a polynomial in either finite-dimensional space
vanishing on that ball vanishes identically.  Multiplying the
basis elements by the bump and the inverse Gram matrix produces
fixed dual moment bumps.  Subtraction removes all moments with
the stated bound on coefficients.
\end{proof}

The preceding inverses reduce the joining problem to estimates for
the moments that they cannot remove. We apply them on the fixed annulus
$\mathcal A=\{z\in\R^n:1<|z|<4\}$. Choose a smooth cutoff
$0\le\chi\le1$, equal to one near its inner sphere and zero near
its outer sphere. For $L$ large put $\chi_L(x)=\chi(x/L)$ on
$\mathcal A_L=\{L<|x|<4L\}$. The following observation explains
why no center-of-mass or angular-momentum limit is required.

\begin{lemma}[Sublinear first moments]\label{lem:end:moments}
Let $h=g-\delta$, $h^{\mathrm b}=g^{\mathrm b}-\delta$, and
$b=K$, $b^{\mathrm b}=K^{\mathrm b}$ on their common coordinate
tail. Put $H=h-h^{\mathrm b}$ and $B=b-b^{\mathrm b}$. Define
the scalar and vector commutators on $\mathcal A_L$ by
\[
 C_L=\mathsf S(\chi_LH)-\chi_L\mathsf S(H),\qquad
 F_L=\mathsf D(\chi_LB)-\chi_L\mathsf D(B),
\]
and their rescaled versions on $\mathcal A$ by
\[
 \widehat C_L(z)=L^2C_L(Lz),\qquad
 \widehat F_L(z)=L^2F_L(Lz).
\]
The flat sources $\mathsf S(H)$ and $\mathsf D(B)$ are integrable.
For every fixed affine scalar function $a(z)$ and every fixed
Euclidean Killing field $Y(z)$,
\[
 \int_{\mathcal A}a\widehat C_L\,dz=o(L^{2-n}),\qquad
 \int_{\mathcal A}Y^i\widehat F_{L,i}\,dz=o(L^{2-n}).
\]
\end{lemma}

\begin{proof}
Expansion of the constraint formulas about $(\delta,0)$ gives
\begin{equation}\label{eq:end:flat-remainders}
 \mathsf S(h)=2\mu+O(r^{-2q-2}),\qquad
 \mathsf D(K)=J+O(r^{-2q-2}).
\end{equation}
The remainder estimates follow, respectively, from
$hD^2h$, $(Dh)^2$, and $K^2$, and from $hDK$ and $Dh\,K$.
The volume densities of $g$ and $\delta$ are uniformly comparable
on the tail.  Hence integrability of $\mu,J$, together with
$2q>n-2$, gives the claim for the original data.  The model is
vacuum with faster decay, so its flat sources are integrable too.

For any integrable tail function $f$, even without an integrable
first absolute moment,
\begin{equation}\label{eq:end:sublinear-moment}
 \int_{R_0<r<4L}r|f|\,dx=o(L).
\end{equation}
Indeed divide by $L$ and split at a fixed large radius $R_1$.
The inner integral divided by $L$ tends to zero, whereas the
outer integral divided by $L$ is at most
$4\int_{r>R_1}|f|$, which tends to zero as $R_1\to\infty$.

Write $Q=H-(\tr_\delta H)\delta$.
For an affine function $a(x)$ define its Green flux by
\[
 \mathcal B_a(R)=\int_{S_R}
 \bigl(a\,\partial_jQ_{ij}\nu^i_\delta
                 -(\partial_i a)Q_{ij}\nu^j_\delta\bigr)\,dA_\delta.
\]
Integration by parts between spheres yields
$\mathcal B_a(R)-\mathcal B_a(R_0)
=\int_{R_0<r<R}a\,\mathsf S(H)$.
Charge matching gives $\mathcal B_1(R)\to0$, and
Equation~\eqref{eq:end:sublinear-moment} gives
$\mathcal B_{x^a}(R)=o(R)$.  For the scalar commutator $C_L$, supported where the cutoff varies,
one consequently has
\begin{equation}\label{eq:end:commutator-moment}
 \int_{L<r<4L}a C_L
 =-\mathcal B_a(L)-\int_{L<r<4L}\chi_La\,\mathsf S(H).
\end{equation}
This is $o(1)$ for $a=1$ and $o(L)$ for $a=x^a$.

For momentum put
$Q_{ij}=(K-K^{\mathrm b})_{ij}
-\tr_\delta(K-K^{\mathrm b})\delta_{ij}$.
For a Euclidean Killing field $Y$, symmetry of $Q$ gives
$Q_{ij}\partial_jY_i=0$.  Its Green flux is thus simply
$\int_{S_R}Y_iQ_{ij}\nu^j_\delta$.
For constant $Y$ its limit is zero by equality of the momenta;
replacing geometric trace reversal by Euclidean trace reversal
changes the sphere flux by $O(R^{n-2-2q})\to0$.
For rotational $Y$ the same Green identity and
Equation~\eqref{eq:end:sublinear-moment} give $o(R)$.
Equation~\eqref{eq:end:commutator-moment} holds for $F_L$ with this flux and pairing.

Under $x=Lz$ the fields are $h(Lz)$ and $LK(Lz)$, and both
flat sources are multiplied by $L^2$.  A constant moment is
therefore multiplied by $L^{2-n}$ and a first moment by
$L^{1-n}$.  The preceding $o(1)$ and $o(L)$ bounds prove the
same $o(L^{2-n})$ bound for all rescaled moments.
\end{proof}

We now apply these two lemmas to the actual joining problem.
Under $x=Lz$, interpolate the scaled fields on $\mathcal A$
between the original and model fields using $\chi$. The errors in
$(2\mu,J)$ relative to the cutoff interpolation of the exact
constraint densities equal $\widehat C_L$ and $\widehat F_L$,
respectively, plus $O(L^{-2q})$ in scaled units. The commutators are
supported in the fixed compact subset where the cutoff varies. By
Lemma~\ref{lem:end:compact-inverses} subtract compact corrections
that cancel these commutators after their moment bumps have
been removed.  The commutators have $W^{1,a}$ norm $O(L^{-q})$,
for any fixed $a>n$, because the available asymptotic derivatives
include those of order three for $g$ and two for $K$.
Sobolev embedding in dimension $n$ therefore gives corrections
of size $O(L^{-q})$ in $C^{2,\alpha}$ for the scaled metric
and in $C^{1,\alpha}$ for the scaled second fundamental form,
where $0<\alpha<1-n/a$.

These are smooth compact corrections; only their first few norms
need be bounded uniformly in $L$.  Their nonlinear contribution
is still $O(L^{-2q})$.  The fixed smooth moment bumps contribute
$o(L^{2-n})$ by Lemma~\ref{lem:end:moments}.  Thus, returning to
physical units, the new data $(\widetilde g_L,\widetilde K_L)$
agree with the original for $r\le L$ and with the model for
$r\ge4L$, and satisfy
\begin{equation}\label{eq:end:deficit}
 \widetilde\mu_L-|\widetilde J_L|_{\widetilde g_L}
 \ge-d_L L^{-n}\,\mathbf 1_{\{L<r<4L\}},\qquad d_L\longrightarrow0.
\end{equation}
To check the unit balls in this assertion, the interpolated
metric, the original metric, and $\delta$ differ by $O(L^{-q})$.
The scaled original $J$ is $O(L^{-q})$.  Transporting a unit
vector between these balls thus changes its pairing with the
interpolated $J$ by $O(L^{-2q})$.  The cutoff is between zero
and one, so the original DEC and the vacuum model give the
claimed sign after precisely the errors already estimated.

\subsection{Changing the vacuum slice and repairing the constraints}

We have joined the original data to a vacuum end with matching
charges.  The only possible failure of the dominant energy condition
has the small bound in Equation~\eqref{eq:end:deficit}.  The next step
removes the momentum by changing the slice entirely within that
vacuum end; it therefore creates no further constraint error.

In the model coordinates, $T=0$ is the graph $b=vy^1$.
Choose a smooth cutoff $\zeta$ equal to one initially and zero
finally, with
\[
 \left|\frac{d\zeta}{d\log |y|}\right|\le\eta_v,
 \qquad v(1+\eta_v)<1.
\]
This can be done by taking a sufficiently long, but fixed,
interval in $\log(|y|/L)$.  Start it outside $|y|=4\gamma L$,
where the joined data are exactly the model.  Replace the graph
by $b=vy^1\zeta(\log(|y|/L))$ there.  Its Euclidean slope is
at most $v(1+\eta_v)<1$; the Schwarzschild coefficients tend
uniformly to their flat values on the transition as $L\to\infty$.
Thus this is a smooth spacelike hypersurface for large $L$,
and its induced data are exactly vacuum throughout the bending
region.  They agree with $b=0$ beyond $|y|=C_vL$, for a fixed
$C_v$.  For $v=0$ no bending is necessary.

Use $y=\operatorname{diag}(\gamma,1,\ldots,1)x$ to parametrize
the whole construction outside the unchanged core.  The resulting
metric and second fundamental form, denoted $(\bar g_L,\bar K_L)$,
satisfy, in this region,
\begin{equation}\label{eq:end:bending-bounds}
 c_v\delta\le\bar g_L\le C_v'\delta,\qquad
 |D\bar g_L|+|\bar K_L|\le C_v'/r.
\end{equation}
On the joining annulus the stronger small perturbation estimates
hold; in the bending region these bounds follow by differentiating
the graph and its normal.  The tensor $\bar K_L$ is compactly
supported, and the only possible DEC deficit is
Equation~\eqref{eq:end:deficit}.

The slice is now static near infinity and its second fundamental
form has compact support.  It remains to repair the annular deficit.
By Equation~\eqref{eq:end:conformal-U}, it is enough to build a
positive function whose negative Laplacian dominates the drift
$|\bar K_L|\,|dU_L|$ and the deficit.  The function must be constant
on the original core, so the boundary expansion keeps its sign.

We construct a nonnegative radial function $f_L(s)$,
$s=|y|/L$, with bounded derivatives on every fixed transition
interval, so that
\begin{equation}\label{eq:end:repair-factor}
 U_L=1+e_LL^{-p}f_L(|y|/L),\qquad e_L\downarrow0,
\end{equation}
repairs that deficit.  Choose a smooth nonnegative $\chi_*$
equal to one on the radial interval containing the full joining
annulus and supported in a slightly larger interval.  Begin
before that interval with $z_L=0$, and solve
\[
 z_L'=C_*z_L+\chi_*,\qquad z_L=-f_L'.
\]
Continue this equation through the bending interval.  The
interval and $C_*$ are fixed for the chosen boost.  Thus $z_L$
and its derivatives are bounded independently of large $L$.
The radial function has $|ds|_{\bar g_L}^2\ge cL^{-2}$,
$|\Delta_{\bar g_L}s|\le CL^{-2}$, and
$|\bar K_L||ds|_{\bar g_L}\le CL^{-2}$ there, by
Equation~\eqref{eq:end:bending-bounds}.  Consequently
\begin{align}
 -\Delta_{\bar g_L}U_L-|\bar K_L|\,|dU_L|_{\bar g_L}
 &=e_LL^{-p}\bigl(z_L'|ds|^2+z_L\Delta s
                             -|\bar K_L|z_L|ds|\bigr)\notag\\
 &\ge c'e_LL^{-n}\chi_*                         \label{eq:end:repair-sign}
\end{align}
when $C_*$ is sufficiently large.  Where $z_L=0$ the factor
is constant, so no condition is needed on the original compact
part.

This construction supplies the required sign through the joining
and bending regions.  To control the ADM charge, we continue it into
the static tail using its radial flux.

There remains a tail extension of $z_L$.  In the exact static
part set $a_L(s)=1+M/(2(Ls)^p)$.  The radial flux factor is
\[
 q_L(s)=s^ka_L(s)^2z_L(s).
\]
The equation just used makes $q_L'>0$ at the beginning of this
static part, after increasing $C_*$ if necessary.  Extend
$q_L$ smoothly so that it stays nondecreasing and is ultimately
a positive constant $q_{L,\infty}$.  For example, continue its
positive derivative over a short interval and then multiply
that derivative by a smooth cutoff that becomes zero.  This
preserves the previous function on an overlap, and all resulting
constants remain bounded independently of large $L$.
Define $z_L=q_L/(s^ka_L^2)$ afterwards and
$f_L(s)=\int_s^\infty z_L(t)\,dt$ everywhere.  Then $f_L\ge0$,
is constant before the transition, and has a finite value there.
The static radial Laplacian satisfies
\[
 \Delta_{\bar g_L}f_L
 =-L^{-2}a_L^{-2n/p}s^{-k}q_L'(s)\le0.
\]
It is zero on the ultimate tail, where $\bar K_L=0$.

Take $e_L\to0$ sufficiently slowly that $d_L/e_L\to0$.
Apply the conformal change
\[
 g_L^*=U_L^{4/p}\bar g_L,\qquad K_L^*=U_L^{2/p}\bar K_L.
\]
Equations~\eqref{eq:end:conformal-U}, \eqref{eq:end:deficit},
and \eqref{eq:end:repair-sign} prove DEC everywhere for all
sufficiently large $L$.  Since $U_L$ is constant on the original
core, the future boundary expansion keeps its weak sign.
Moreover $U_L\ge1$, so this repair increases the metric as a
quadratic form.

On the static tail, $f_L(s)=q_{L,\infty}p^{-1}s^{-p}
+O(L^{-p}s^{-2p})$.  It follows that the coefficient of $|y|^{-p}$
in $U_L-1$ is $e_Lq_{L,\infty}/p$.  Computed in the
asymptotically Euclidean rest coordinates $y$, the energy is
\begin{equation}\label{eq:end:repair-energy}
 E_L^*=M+\frac{2e_Lq_{L,\infty}}p=M+o(1),\qquad P_L^*=0.
\end{equation}
One can also see the exact normalization without an expansion:
on this tail $a_LU_L$ is a radial Euclidean harmonic function
with limit one, because $a_L^{4/p}\delta$ is scalar flat and
$U_L$ is harmonic for that metric.  Thus the ultimate metric
is exactly an isotropic Schwarzschild metric with the mass in
Equation~\eqref{eq:end:repair-energy}.  The new constraints
are integrable, since the repair source is compactly supported
and the ultimate end is vacuum.

\subsection{Comparing every full enclosing cut}

We have obtained zero momentum, energy tending to $m$, and the
dominant energy condition.  The numerical transfer also requires a
lower bound for the infimum over \emph{all} full cuts.  The metric
need not approach the original metric uniformly on
the whole end during the change of slice.  The following
comparison supplies what is needed for all enclosing areas.
In the original $x$ coordinates the metrics $g_L^*$ have a
uniform lower bound $c_0\delta$ throughout the region $r\ge L$,
where $c_0>0$ depends on the fixed boost.  Choose a fixed
$0<\lambda<1$ with $\lambda^2<c_0/2$.

Let $h_L:[0,\infty)\to[0,\infty)$ be smooth and increasing,
equal to $r$ for $r\le\sqrt L$, with $\lambda\le h_L'\le1$,
and with $h_L'=\lambda$ for $r\ge2\sqrt L$.
Set $\Phi_L(r\theta)=h_L(r)\theta$ on the end, and extend by
the identity on the core.  The derivative eigenvalues in the
Euclidean radial and tangential directions are $h_L'$ and
$h_L/r$, respectively.  They are at most one; for $r\ge L$
they are at most $\lambda+O(L^{-1/2})$, since
$h_L(r)=\lambda r+O(\sqrt L)$.  Also $h_L\to\infty$, so
$\Phi_L$ is a proper diffeomorphism of the entire exterior
onto itself.

For $r\le\sqrt L$ the map is the identity and $g_L^*\ge g$.
For $\sqrt L\le r\le L$, both $g$ at $x$ and $g$ at
$\Phi_L(x)$ equal $\delta+o(1)$ uniformly, and $g_L^*$ is
a conformal multiple, at least one, of $g$.
The derivative bound therefore gives
$\Phi_L^*g\le(1+o(1))g_L^*$ there.  For $r\ge L$ one instead
has
\[
 \Phi_L^*g\le(\lambda^2+o(1))\delta\le g_L^*.
\]
Combining the regions proves
\begin{equation}\label{eq:end:compression}
 \Phi_L^*g\le(1+o(1))g_L^*.
\end{equation}

If $\Gamma=\partial D$ is any full enclosing cut for $g_L^*$,
then $\Phi_L(D)$ is connected, closed, has interior in
$\operatorname{int}\Omega$, and contains the entire distant
end.  Its full intrinsic boundary is $\Phi_L(\Gamma)$,
including every component and every portion on $S$.
Restricting Equation~\eqref{eq:end:compression} to each
$k$-dimensional tangent plane and taking determinants gives
\begin{equation}\label{eq:end:all-cut-comparison}
 \Area_{g_L^*}(\Gamma)
 \ge(1+o(1))^{-k/2}\Area_g(\Phi_L(\Gamma))
 \ge(1-o(1))a.
\end{equation}
Taking an infimum uses no existence or regularity of a minimizing
cut.

\subsection{A strict scalar perturbation}

The use of a spacetime Poisson equation for conformal strictification
has a direct precedent in Jaracz's three-dimensional construction
\cite[Theorem~1.1 and Section~3.1, preprint version~1]{Jaracz2025StrictDEC}.
The conformal scaling below is its dimension-dependent counterpart.
Here a negative inner normal derivative also makes trapping strict;
the smoothed drift, prescribed decay, and charge and full-cut estimates
are established for the present exterior class.

At this stage the inequalities may still be non-strict.  We solve a
scalar equation with positive source and negative inner normal
derivative.  The source gives a quantitative strict energy margin;
the normal derivative makes the future expansion strictly negative.
We state the equation on the original decay class as well as on a
static tail, so the resulting scalar perturbation is available
independently of the preceding replacement.

\begin{lemma}[Strictification]\label{lem:end:strict}
Let $\Omega$ be a smooth connected exterior, complete with its
nonempty compact smooth boundary $S$ included, and with exactly one
Euclidean coordinate end outside a compact set.  Let $(g,K)$ be
smooth up to $S$, with $g-\delta=O_6(r^{-q})$ and
$K=O_5(r^{-1-q})$ for some $q>(n-2)/2$.  Choose $0<\beta<\min(1,q)$ and a smooth
nonnegative function $B\ge|K|_g$ with
$B=O_5(r^{-1-q})$.  Then the problem
\begin{equation}\label{eq:end:strict-pde}
 -\Delta_g\phi
 =B\sqrt{|d\phi|_g^2+\varrho^{-2k}}+\varrho^{-n-\beta},
 \qquad \partial_\nu\phi=-1\text{ on }S,
 \qquad \phi\longrightarrow0\text{ at infinity}
\end{equation}
has a smooth positive solution, unique among decaying smooth
solutions.  It satisfies
\begin{equation}\label{eq:end:strict-estimates}
\begin{gathered}
 \phi=O_6(r^{2-n}),\qquad
 \frac{-\Delta_g\phi}{r^{-n-\beta}}\longrightarrow1,\\
 \mathfrak f_\phi:=\lim_{R\to\infty}
          \int_{S_R}\partial_{\nu_g}\phi\,dA_g\\
 =-\Area_g(S)-\int_\Omega
       \bigl(B\sqrt{|d\phi|^2+\varrho^{-2k}}
                    +\varrho^{-n-\beta}\bigr)\,dV_g<0.
\end{gathered}
\end{equation}
If $K$ is compactly supported, $B$ may be chosen compactly
supported.  If in addition $g-\delta=O_d(r^{2-n})$ for all $d$,
the solution has this decay through every derivative order.

Suppose also that the ADM limits of $(g,K)$ exist and are finite,
that the data satisfy DEC, and that $\theta_+(S)\le0$ on every
boundary component.  For all sufficiently small $\delta>0$,
\[
 g_\delta=e^{2\delta\phi}g,\qquad
 K_\delta=e^{\delta\phi}K
\]
satisfy
\begin{equation}\label{eq:end:strict-margin}
 \mu_\delta-|J_\delta|_{g_\delta}\ge c_\delta\varrho^{-n-\beta},
 \qquad \theta_{+,\delta}<0,
 \qquad E_\delta=E-\frac{\delta}{\omega}\mathfrak f_\phi,
 \qquad P_\delta=P.
\end{equation}
Here $g_\delta\ge g$, and their ratio tends uniformly to one
as $\delta\downarrow0$.
\end{lemma}

\begin{proof}
\medskip\noindent\emph{1. Continuation on a truncation.}
Write $a=\varrho^{-k}$ and $b=\varrho^{-n-\beta}$.  The
positive regularizing term $a^2$ makes the right side smooth
as a function of $d\phi$, with first gradient derivative of
norm at most $B$.  Choose a large coordinate sphere enclosing
$S$ and a compact core.  On a larger truncation $\Omega_R$
impose $\phi=0$ on $S_R$, and multiply $B$ in the equation by
$t\in[0,1]$.  The inner Neumann and outer Dirichlet boundary
components are disjoint smooth hypersurfaces.

At $t=0$ this mixed Laplace problem is solvable.  Indeed its
weak formulation is coercive on $H^1$ functions of zero outer
trace by Poincare's inequality; the Neumann datum is a bounded
boundary functional by the trace theorem.  Lax--Milgram gives
the weak solution and local boundary regularity gives a smooth
solution.  The linearization at any solution for $0\le t\le1$
is
\[
 -\Delta_g u-tB\frac{\langle d\phi,du\rangle_g}
                     {\sqrt{|d\phi|^2+a^2}},
 \qquad \partial_\nu u=0\text{ on }S,\quad u=0\text{ on }S_R.
\]
It has a bounded drift, is uniformly elliptic, and has no
zeroth order term.  Its homogeneous kernel is zero by the
strong maximum principle and the Hopf boundary lemma: a
nonzero maximum or minimum cannot occur in the interior or
on a zero-Neumann component \cite[Chapter~3]{GilbargTrudinger}.
The same conclusion follows
for weak limits with merely bounded drift.  The drift is a
compact perturbation of the mixed Laplace isomorphism in
Hölder spaces, so its index is zero and this injectivity gives
surjectivity.  The implicit function theorem gives openness
of the continuation set.

\medskip\noindent\emph{2. Bounds independent of the truncation.}
We detail the estimates giving closedness and exhaustion.
All solutions are nonnegative.  An interior negative minimum
contradicts $-\Delta\phi>0$.  At a minimum on $S$ the Hopf
lemma requires a negative outward derivative, whereas the
outward derivative of $\Omega_R$ is $-\partial_\nu\phi=1$.
The outer boundary value is zero.

For $0<\alpha<\min(\beta,q)$ and sufficiently large $r_0$,
the function
\[
 W(r)=r^{-p}\bigl(1-r^{-\alpha}\bigr),\qquad r\ge r_0,
\]
is positive and satisfies
\begin{equation}\label{eq:end:strict-barrier}
 -\Delta_gW-B|dW|_g\ge c r^{-n-\alpha}.
\end{equation}
In fact $-\Delta_\delta W=\alpha(p+\alpha)r^{-n-\alpha}$,
the metric error is $O(r^{-n-q})$, and
$B|dW|=O(r^{-n-q})$.  Since $b+Ba=O(r^{-n-\beta})$,
a sufficiently large multiple $C(A+1)W$, where
$A=\sup_{S_{r_0}}\phi$, is a supersolution outside $r_0$,
dominates $\phi$ on $S_{r_0}$, and is nonnegative on $S_R$.
Here we used
$\sqrt{|d\phi|^2+a^2}\le|d\phi|+a$ and the fact that this
gradient dependence is Lipschitz.  Subtracting a supersolution
therefore gives a scalar linear inequality with bounded drift,
to which comparison applies.  Hence
\begin{equation}\label{eq:end:strict-tail-bound}
 0\le\phi(x)\le C(1+A)r^{-p}\quad(r_0\le r\le R),
\end{equation}
uniformly in $R$ and $t$.

We claim that the suprema of all these solutions are uniformly
bounded.  Otherwise divide a sequence by its suprema $M_i\to
\infty$.  Estimate~\eqref{eq:end:strict-tail-bound} ensures
that the normalized functions have a positive supremum on
a fixed compact region.  Local $W^{2,a_0}$ estimates, for any
fixed $a_0>n$, apply uniformly, including the inner Neumann
boundary: the equation has right side bounded by a fixed
multiple of $|d\phi|+1$, and first derivative interpolation
absorbs the gradient term in the second derivative estimate.
The same argument applies after division by $M_i$.
Compactness gives a nonnegative, nonzero limit $u$ with
$u\to0$ at infinity and
\[
 -\Delta_gu=t_*B|du|,\qquad \partial_\nu u=0.
\]
If the outer radii instead have a finite limit, the limit has
zero data at that limiting sphere.  In either case its positive
maximum is attained.  The strong maximum principle, or the
Hopf lemma at $S$, contradicts that maximum.  These principles
apply because $B|du|$ equals a bounded measurable drift paired
with $du$, defining that drift to be zero where $du=0$.
This proves the uniform supremum estimate.

\medskip\noindent\emph{3. Exhaustion and uniqueness.}
The $W^{2,a_0}$ estimates now give uniform $C^{1,\gamma}$
bounds on fixed patches.  Their right sides are then Hölder
continuous, and the scalar Schauder estimates give
$C^{2,\gamma}$ bounds.  Differentiation and the strictly
positive regularizer $a$ give all higher local bounds.
For each fixed truncation these estimates close continuation
to $t=1$.  Letting $R\to\infty$ gives a smooth solution of
Equation~\eqref{eq:end:strict-pde}; the tail bound gives the
specified limit.  It is strictly positive by the same minimum
principle.  The difference of two decaying solutions has zero
Neumann data and solves a homogeneous equation with bounded
drift, by integrating the gradient derivative of the square
root along the segment between their gradients.  Its positive
and negative extrema are excluded as above, proving uniqueness.

\medskip\noindent\emph{4. End decay and flux.}
On dyadic end annuli rescale $x=Lz$ and
$\Phi(z)=L^p\phi(Lz)$.  Estimate~\eqref{eq:end:strict-tail-bound}
bounds $\Phi$.  The rescaled equation has uniformly controlled
metric coefficients and the form
\[
 -\Delta_{g(Lz)}\Phi
 =O_{5}(L^{-q})\sqrt{|d\Phi|_{g(Lz)}^2+|z|^{-2k}}
                       +L^{-\beta}|z|^{-n-\beta},
\]
where the coefficient denoted $O_5(L^{-q})$ is nonnegative
and smooth.  The same interpolation, scalar estimates and
differentiation give uniform bounds for six derivatives of
$\Phi$.  In the compactly supported $B$ case the equation is
just a Poisson equation on the distant end, and all available
metric derivative orders give corresponding derivative
bounds.  This proves the asserted decay.  It also gives
$B\sqrt{|d\phi|^2+r^{-2k}}=O(r^{-n-q})$, so $\beta<q$
proves the Laplacian ratio in
Equation~\eqref{eq:end:strict-estimates}.

The source is integrable.  Integrating the equation on
$\Omega_R$, remembering that its inner outward normal is
$-\nu$, gives
\[
 \int_{S_R}\partial_{\nu_g}\phi\,dA_g
 =-\Area_g(S)-\int_{\Omega_R}
              \bigl(B\sqrt{|d\phi|^2+a^2}+b\bigr)\,dV_g.
\]
This proves the claimed finite negative flux.

\medskip\noindent\emph{5. Strict inequalities and charges.}
Finally $|d\phi|^2\le Cb$ globally: on the end the quotient
has order $r^{-(n-2-\beta)}$, which is bounded because
$\beta<1\le n-2$, and on the core $b$ is positive.
Equation~\eqref{eq:end:strict-pde} implies, for every unit-ball
vector $v$,
$-\Delta\phi+K(\nabla\phi,v)\ge b$.
Use Equation~\eqref{eq:end:conformal} with $s=\delta\phi$;
for sufficiently small fixed $\delta>0$ the negative square
term is at most half the positive $k\delta b$ term.
This proves the strict DEC margin, and
$e^{\delta\phi}\theta_{+,\delta}=\theta_+-k\delta<0$.
The conformal energy vector changes by
$-2k\delta\,d\phi+o(r^{1-n})$.  The geometric and Euclidean
fluxes of $d\phi$ agree in the limit, since their difference
is $O(r^{-q})$ after integration.  Division by $2k\omega$
gives the energy formula in
Equation~\eqref{eq:end:strict-margin} exactly.  The momentum
trace reversal is multiplied by $e^{\delta\phi}$; its flux
change is $O(r^{-q})\to0$.  Positivity and boundedness of
$\phi$ give the last metric statements.
\end{proof}

\paragraph{Infinitesimal strictification.}
For any datum in Lemma~\ref{lem:end:strict}, the same construction
gives a quantitative infinitesimal statement at points where the
dominant energy condition is saturated. Write
$\mathfrak C_\delta(v)=2\bigl(\mu_\delta+
J_\delta(e^{-\delta\phi}v)\bigr)$.  At an active ray,
meaning $|v|_g\le1$ and $\mu+J(v)=0$, differentiation of
Equation~\eqref{eq:end:conformal} gives
\begin{equation}\label{eq:end:strict-direction}
 \mathfrak C'_0(v)=2k\bigl(-\Delta\phi+K(\nabla\phi,v)\bigr)
       \ge2k\varrho^{-n-\beta}.
\end{equation}
Its quotient by $r^{-n-\beta}$ tends to $2k$ along any active
rays escaping to infinity: the two terms involving $B$ and $K$
are $O(r^{-n-q})$.  At a marginal boundary the expansion
derivative of the same direction is exactly $-k$.

\begin{proof}[Proof of Proposition~\ref{prop:end:reduction}]
Apply Lemma~\ref{lem:end:strict} to the repaired data
$(g_L^*,K_L^*)$, choosing $B$ compactly supported and choosing
$\delta_L\downarrow0$ so that
$\delta_L\|\phi_L\|_\infty\to0$ and
$\delta_L|\mathfrak f_{\phi_L}|\to0$.
There is no need for constants in this solve to be uniform in
$L$: this choice is made after the entire replacement at that
$L$.  Its scalar curvature has the required decay, since the
ultimate initial end is scalar flat and
\[
 R_{e^{2\delta_L\phi_L}g_L^*}
 =e^{-2\delta_L\phi_L}
       \bigl(2k\delta_L r^{-n-\beta}
                         -kp\delta_L^2|d\phi_L|^2\bigr)
\]
there.  Compact support of $K_L^*$ is preserved, and all metric
derivative orders have the stated asymptotics.  The strict
margin and expansion are those of
Lemma~\ref{lem:end:strict}.  The core remains smooth and the
tail is uniformly Euclidean comparable, so completeness with
$S$ included is preserved.  The constraints are integrable
on the core and have integrable decay on the tail.

By Equation~\eqref{eq:end:repair-energy} and the choice of
$\delta_L$, the final energy tends to $M=m>0$, and the final
momentum is zero. Strictification only increases the metric,
so Equation~\eqref{eq:end:compression} and the full-cut
comparison remain valid; in particular, every resulting full-cut
infimum is positive. Pass to a sequence along which the comparison
errors are bounded by numbers $\epsilon_j\downarrow0$.
These data satisfy Definition~\ref{def:prepared} and every assertion
in Equation~\eqref{eq:end:reduction-limits}.
\end{proof}

\subsection{A positive shell with fixed momentum}

The reduction above uses a fixed original future-timelike ADM vector.
The following independent conformal change applies without a timelike
assumption: it raises the energy by a prescribed amount, fixes the
momentum, and increases every full-cut area. These are the properties
used in the proof of Corollary~\ref{cor:automatic-timelikeness} in
Section~\ref{sec:comparison}.

\begin{lemma}[A positive shell in the strong-decay class]
\label{lem:end:positive-shell}
Assume all hypotheses of Theorem~\ref{thm:main} except possibly
$E>|P|_\delta$. For every $\lambda>0$ there is a smooth function
$U_\lambda\ge1$, constant near $S$, such that
\[
 g_\lambda=U_\lambda^{4/p}g,\qquad
 K_\lambda=U_\lambda^{2/p}K
\]
are smooth complete one-ended data with the same $O_6/O_5$ decay and $q$,
integrable constraint densities, the dominant energy condition and weakly
future trapped boundary, and satisfy
\[
 E_\lambda=E+2\lambda,\qquad P_\lambda=P,\qquad
 A_{\min,g_\lambda}(S)\ge A_{\min,g}(S)>0.
\]
\end{lemma}

\begin{proof}
For $U>0$ put $a=U^{2/p}$ and $(g_U,K_U)=(a^2g,aK)$.
Substitution of $s=2p^{-1}\log U$ into
Equation~\eqref{eq:end:conformal} gives, as scalar and covector identities,
\[
 a^2\mu_U=\mu-\frac{2k}{pU}\Delta_gU,\qquad
 aJ_U=J+\frac{2k}{pU}K(\nabla U,\cdot).
\]
Indeed, the scalar terms combine as
$-k\Delta s-\frac{kp}{2}|ds|^2=-\frac{2k}{pU}\Delta U$.
Taking the new-metric covector norm introduces another factor $a^{-1}$,
so
\begin{equation}\label{eq:end:shell-dec}
 a^2(\mu_U-|J_U|_{g_U})
 \ge \mu-|J|_g+\frac{2k}{pU}
       \bigl(-\Delta_gU-|K|_g|dU|_g\bigr).
\end{equation}
The boundary identity is
\[
 a\theta_{+,U}=\theta_++\frac{2k}{pU}\partial_\nu U.
\]
Thus a factor constant near $S$ preserves its expansion sign, and
$-\Delta_gU\ge |K|_g|dU|_g$ preserves the dominant energy condition.

Choose $0<\xi<\min(q,1)$. In a sufficiently distant part of the
coordinate end set
\[
 T(r)=r^{-p}-r^{-p-\xi}.
\]
There $T>0$ and $T'<0$. Since $p=n-2$,
$\Delta_\delta r^{-\alpha}=\alpha(\alpha+2-n)r^{-\alpha-2}$.
The original metric and tensor estimates therefore give
\[
 -\Delta_gT-|K|_g|dT|_g
 =\xi(p+\xi)r^{-n-\xi}+O(r^{-n-q})>0
\]
after increasing the radius. Choose a nondecreasing smooth switch
$\beta$ from zero to one in this region, constant near its endpoints,
and put
\[
 G(r)=-\int_r^\infty\beta(t)T'(t)\,dt.
\]
Extend $G$ constantly through the interior. Then $G\ge0$, it is constant
near $S$, and $G=T$ sufficiently far out. Because $dG=\beta\,dT$,
\[
 -\Delta_gG-|K|_g|dG|_g
 =\beta(-\Delta_gT-|K|_g|dT|_g)-\beta'T'|dr|_g^2\ge0.
\]
Use $U_\lambda=1+\lambda G$. Equation~\eqref{eq:end:shell-dec} and
the boundary formula prove the required signs for every fixed
$\lambda>0$, without a smallness condition.

The factor is bounded for each fixed $\lambda$ and is at least one,
so completeness with $S$ included is preserved. The manifold and
its unique coordinate end are unchanged. On the tail $G=O_j(r^{-p})$
at every fixed derivative order. Since $q<p$, the added metric tail
is faster than the original decay, and
\[
 g_\lambda-\delta=O_6(r^{-q}),\qquad
 K_\lambda=O_5(r^{-1-q}).
\]
Moreover,
\[
 \Delta_gU_\lambda=O(r^{-n-\xi})+O(r^{-n-q}),\qquad
 |K|_g|dU_\lambda|_g=O(r^{-n-q}).
\]
These terms are integrable against $dV_g\asymp r^{n-1}dr\,d\omega$.
The scalar and covector identities above, and the bounded conformal
factors for fixed $\lambda$, show that the new constraint densities
are integrable against $dV_{g_\lambda}$.

For the charges write $f_\lambda=U_\lambda^{4/p}$. Then
\[
 f_\lambda=1+\frac{4\lambda}{p}G+O_1(r^{-2p}).
\]
In the metric flux of Equation~\eqref{eq:setting:energy}, the leading
change is $-k\partial_i(f_\lambda-1)$. Terms involving $g-\delta$ or
its first derivatives have integrated flux $O(R^{-q})$, and the nonlinear
remainder has flux $O(R^{-p})$. Since
\[
 \lim_{R\to\infty}\int_{S_R}\partial_rG\,dA_\delta
 =-p\omega_{n-1},
\]
the current normalization gives
\[
 E_\lambda-E
 =-\frac{2\lambda}{p\omega_{n-1}}
   \lim_{R\to\infty}\int_{S_R}\partial_rG\,dA_\delta
 =2\lambda.
\]
If $\pi=K-(\tr_gK)g$, its transformed tensor is exactly
$\pi_\lambda=U_\lambda^{2/p}\pi$. The additional momentum flux in
Equation~\eqref{eq:setting:momentum} is $O(R^{-q})$, because
$U_\lambda^{2/p}-1=O(r^{-p})$ and $\pi=O(r^{-1-q})$.
Thus $P_\lambda=P$, and the new ADM limits are finite.

Finally, the admissible domains in Definition~\ref{def:fullcuts} are
unchanged. For their entire intrinsic boundaries,
\[
 \Area_{g_\lambda}(\partial D)
 =\int_{\partial D}U_\lambda^{2k/p}\,dA_g
 \ge \Area_g(\partial D).
\]
This comparison includes every component and boundary-coincident portion.
Taking the infimum proves the claimed full-cut area inequality.
\end{proof}
